\documentclass[12pt,a4paper]{article}

\usepackage[margin=1in]{geometry}
\usepackage{setspace}
\usepackage{hyperref}
\usepackage{enumitem}
\usepackage{longtable}

\begin{document}

%------------------ Title Page ------------------%
\begin{titlepage}
\thispagestyle{empty}
\begin{center}
\vspace*{3cm}
{\LARGE \textbf{Scientific Calculator DevOps Pipeline}\par}

\vspace{2cm}

\textbf{Abstract:}\\[0.5cm]
This project demonstrates a complete end-to-end CI/CD lifecycle for a Python-based scientific calculator.
By integrating Jenkins for orchestration, Docker for containerization, and Ansible for automated deployment,
the pipeline ensures that every code change is rigorously tested and seamlessly delivered to a production-ready environment.

\vfill

\begin{flushright}
\textbf{Project By:}\\[0.5cm]
\textbf{Name:} Kartikey Dubey\\
\textbf{Roll Number:} MT2025061
\end{flushright}

\end{center}
\end{titlepage}

\setstretch{1.2}

%------------------ Main Report ------------------%

\section*{Project Report: Scientific Calculator DevOps Pipeline}

\section{Basic Ideas: The What and Why of DevOps}

\textbf{What is DevOps?}\\
DevOps is a set of practices that combines software development (Dev) and IT operations (Ops). It aims to shorten the systems development life cycle and provide continuous delivery with high software quality.

\vspace{0.5em}
\textbf{Why DevOps?}
\begin{itemize}[noitemsep]
  \item \textbf{Speed:} Deliver features and fixes faster.
  \item \textbf{Reliability:} Ensure the quality of application updates and infrastructure changes.
  \item \textbf{Scalability:} Manage infrastructure and development processes at scale.
  \item \textbf{Collaboration:} Build more effective teams under a model that emphasizes ownership and accountability.
\end{itemize}

\section{Tools Used in This Project}

\begin{longtable}{|p{3cm}|p{3cm}|p{7cm}|}
\hline
\textbf{Category} & \textbf{Tool} & \textbf{Purpose} \\
\hline
Source Control & Git / GitHub & Versioning code and hosting the remote repository. \\
\hline
CI/CD Orchestration & Jenkins & Automating the build, test, and deployment stages. \\
\hline
Containerization & Docker & Packaging the application and its dependencies into a portable image. \\
\hline
Artifact Repository & Docker Hub & Storing and sharing built Docker images. \\
\hline
Config Management & Ansible & Automating the deployment and configuration of the container on the host. \\
\hline
Programming Language & Python & Developing the scientific calculator application logic. \\
\hline
\end{longtable}

\section{Step-by-Step Instructions to Reproduce}

\subsection{Phase 1: Source Control with Git}

Source control is the foundation of DevOps. It allows multiple developers to work on the same codebase simultaneously without overwriting each other's work.

\textbf{What is Git?} Git is a distributed version control system that tracks changes in source code during software development.

\vspace{0.5em}
\textbf{Setup \& Commands:}
\begin{enumerate}[noitemsep]
  \item Initialize the repository: \texttt{git init}
  \item Add files: \texttt{git add .}
  \item Commit changes: \texttt{git commit -m "Initial commit"}
  \item Rename branch: \texttt{git branch -M main}
  \item Set remote: \texttt{git remote add origin https://github.com/kartikeyblaze/SPE-mini-project-calc.git}
  \item Push to GitHub: \texttt{git push -u origin main}
\end{enumerate}

% SPACE FOR SCREENSHOT: Git push success in terminal showing remote upload

\subsection{Phase 2: Jenkins Orchestration}

\textbf{Intro:} Jenkins is an open-source automation server that helps automate the parts of software development related to building, testing, and deploying.

\textbf{Setup \& Configuration:}
\begin{enumerate}[noitemsep]
  \item \textbf{Expose Jenkins with ngrok:} Run \texttt{ngrok http 8080} to get a public URL for your local Jenkins. % Terminal showing ngrok forwarding URL
  \item \textbf{GitHub Webhook:} In GitHub Settings $\rightarrow$ Webhooks, add your ngrok URL followed by \texttt{/github-webhook/}. Set Content type to \texttt{application/json}. % GitHub Webhook configuration screenshot
  \item \textbf{Docker Hub Credentials:} In Jenkins $\rightarrow$ Credentials $\rightarrow$ System $\rightarrow$ Global credentials, add a ``Username with password'' entry with ID \texttt{DockerHubCred}. % Jenkins Credentials screenshot
  \item \textbf{Mail Notification:} Configure the SMTP server in Jenkins $\rightarrow$ Manage Jenkins $\rightarrow$ System. Set the SMTP server (e.g., \texttt{smtp.gmail.com}) and authentication, use port 587 for TLS, and 465 if using only SSL. % Jenkins Email Notification screenshot
\end{enumerate}

\subsection{Phase 3: Jenkinsfile Configuration}

The \texttt{Jenkinsfile} defines the entire CI/CD pipeline as code.

\begin{verbatim}
pipeline {
    agent any

    environment {
        DOCKER_HUB_USER = 'kartikeyblaze'
        IMAGE_NAME = 'scientific-calculator'
        IMAGE_TAG = "${env.BUILD_ID}"
        DOCKER_CREDENTIALS_ID = 'DockerHubCred'
        EMAIL = 'kartikeyblaze@gmail.com'
    }

    triggers { 
      githubPush() 
    }

    stages {
        stage('Checkout') {
            steps {
                checkout scm
            }
        }

        stage('Unit Testing') {
            steps {
                sh 'python3 -m unittest test_calculator.py'
            }
        }

        stage('Build Docker Image') {
            steps {
                sh "docker build -t ${DOCKER_HUB_USER}/${IMAGE_NAME}:${IMAGE_TAG} ."
                sh "docker tag ${DOCKER_HUB_USER}/${IMAGE_NAME}:${IMAGE_TAG} ${DOCKER_HUB_USER}/${IMAGE_NAME}:latest"
            }
        }

        stage('Push to Docker Hub') {
            steps {
                script {
                    withCredentials([usernamePassword(credentialsId: DOCKER_CREDENTIALS_ID, passwordVariable: 'DOCKER_PASSWORD', usernameVariable: 'DOCKER_USERNAME')]) {
                        sh 'echo "$DOCKER_PASSWORD" | docker login -u "$DOCKER_USERNAME" --password-stdin'
                        sh "docker push ${DOCKER_HUB_USER}/${IMAGE_NAME}:${IMAGE_TAG}"
                        sh "docker push ${DOCKER_HUB_USER}/${IMAGE_NAME}:latest"
                    }
                }
            }
        }

        stage('Run Ansible Playbook') {
            steps {
                ansiblePlaybook(
                    playbook: 'deploy.yml',
                    inventory: 'inventory'
                )
            }
        }
    }

    post {
        always {
            script {
                def status = currentBuild.result ?: 'SUCCESS'
                mail to: EMAIL,
                     subject: "Build ${status}: ${env.JOB_NAME} #${env.BUILD_NUMBER}",
                     body: "Build ${status} for ${env.JOB_NAME} #${env.BUILD_NUMBER}.\nView more: ${env.BUILD_URL}"
            }
        }
    }
}
\end{verbatim}

\textbf{Block-wise Explanation:}
\begin{itemize}[noitemsep]
  \item \texttt{pipeline \{ agent any \}}: Directs Jenkins to run this pipeline on any available agent.
  \item \texttt{environment \{ ... \}}: Defines global variables like \texttt{DOCKER_HUB_USER} and \texttt{IMAGE_NAME} used throughout the stages.
  \item \texttt{triggers \{ githubPush() \}}: Tells Jenkins to start the build whenever a push event is received from the GitHub webhook.
  \item \texttt{stage('Unit Testing')}: Runs \texttt{unittest} on \texttt{test\_calculator.py}. If tests fail, the pipeline stops.
  \item \texttt{stage('Build Docker Image')}: Builds the image using the local \texttt{Dockerfile} and tags it with the build ID and \texttt{latest}.
  \item \texttt{stage('Push to Docker Hub')}: Securely logs into Docker Hub using \texttt{withCredentials} and pushes the images.
  \item \texttt{stage('Run Ansible Playbook')}: Triggers the deployment phase using the Ansible plugin.
  \item \texttt{post \{ always \{ ... \} \}}: Sends an email notification regardless of whether the build succeeded or failed.
\end{itemize}

% SPACE FOR Jenkins dashboard and pipeline stage screenshots

\subsection{Phase 4: Containerization with Docker}

\textbf{Intro:} Docker allows you to package an application with all of its requirements (like Python libraries) into a container that runs the same on any machine.

\textbf{Dockerfile Explanation:}
\begin{verbatim}
FROM python:3.9-slim  # Uses a lightweight Python 3.9 image
WORKDIR /app          # Sets the internal directory for the app
COPY Calculator.py .  # Copies the script into the container
CMD ["python", "Calculator.py"]  # Command to run when the container starts
\end{verbatim}

\textbf{Commands:}
\begin{enumerate}[noitemsep]
  \item Building: \texttt{docker build -t scientific-calculator .}
  \item Running: \texttt{docker run -it scientific-calculator}
\end{enumerate}

% SPACE FOR Docker build / run and Docker Hub screenshots

\subsection{Phase 5: Deployment with Ansible}

\textbf{Intro:} Ansible is an IT automation tool that automates cloud provisioning, configuration management, and application deployment.

\textbf{Inventory File (\texttt{inventory}):}\\
Defines the target hosts. In this project, it points to \texttt{localhost} for deployment.
\begin{verbatim}
[localhost]
localhost ansible_python_interpreter=/usr/bin/python3
\end{verbatim}

\textbf{Playbook Explanation (\texttt{deploy.yml}):}
\begin{verbatim}
---
- name: Pull Docker Image from Docker Hub and start the container
  hosts: localhost
  connection: local
  gather_facts: no
  tasks:
    - name: Pull Docker Image
      docker_image:
        name: "kartikeyblaze/scientific-calculator:latest"
        source: pull

    - name: Remove existing container if it exists
      docker_container:
        name: "calculator"
        state: absent

    - name: Start Docker service
      service:
        name: docker
        state: started

    - name: Running container
      shell: docker run -it -d --name calculator kartikeyblaze/scientific-calculator:latest
\end{verbatim}

\textbf{Key Tasks:}
\begin{itemize}[noitemsep]
  \item \texttt{docker\_image}: Pulls the latest image from the Docker Hub repo \texttt{kartikeyblaze/scientific-calculator}.
  \item \texttt{docker\_container}: Ensures any old version of the container is removed (\texttt{state: absent}) to avoid naming conflicts.
  \item \texttt{shell}: Runs the new container in detached mode (\texttt{-d}) using the \texttt{latest} image.
\end{itemize}

% SPACE FOR ansible-playbook and docker ps screenshots

\subsection*{Email Notifications}

These are sent regardless of the build status (SUCCESS or FAILURE) along with a build number and link for console output of the specified build.

% SPACE FOR email notification screenshots

\section*{Conclusion \& Future Scope}

This project successfully implemented a standard DevOps lifecycle. The automated pipeline reduces human intervention, minimizes deployment errors, and ensures that only code passing all unit tests reaches production.

\end{document}
